% auto-generated figure: class-competition (source: wf)
\begin{figure}[H]
\centering
\begin{tikzpicture}[node distance=10mm]
  \UmlClass{comp}{Competition}{%
    - name : String\\
    - type : CompetitionType\\
    - status : CompetitionStatus}{%
    + isCompleted()\\
    + standings()}
  \UmlClass[below left=16mm and -22mm of comp]{team}{CompetitionTeam}{%
    - name : String\\
    - shortName : String\\
    - crestUrl : String\\
    - sport : Sport}{}
  \UmlClass[below right=16mm and -22mm of comp]{fix}{CompetitionFixture}{%
    - homeTeamId : Long\\
    - awayTeamId : Long\\
    - group : String [0..1]\\
    - matchday : int\\
    - round : String\\
    - homeScore : Integer\\
    - awayScore : Integer}{%
    + isPlayed()}

  % compositions (filled diamond at the Competition / whole end)
  \draw[{Diamond[fill=dLine,length=3mm,width=2mm]}-,draw=dLine,thick]
        (comp.south) -- node[lbl,pos=0.85,left]{1..*} (team.north);
  \draw[{Diamond[fill=dLine,length=3mm,width=2mm]}-,draw=dLine,thick]
        (comp.south) -- node[lbl,pos=0.85,right]{1..* \{ordered\}} (fix.north);

  \node[lbl,align=left,text width=52mm,font=\scriptsize,
        below=4mm of fix.south,anchor=north]
       {\itshape round distinguishes the Champions League league
        phase (e.g.\ ``League Phase'') from the knockout rounds
        (e.g.\ ``Round of 16'', ``Final'').};
\end{tikzpicture}
\caption{Class model of the competitions domain: a Competition composes its participating CompetitionTeam rows and its ordered CompetitionFixture rows, whose round label separates the Champions League league phase from the knockout rounds.}
\label{fig-class-competition}
\end{figure}
